%==============================================================
%  lecons/algebre/spectral.tex — Théorème spectral
%==============================================================

\lecontitle{Théorème spectral}{Algèbre linéaire — Endomorphismes autoadjoints}

%=============================================================
%  § 1 — ÉNONCÉ
%=============================================================
\secnum{navyblue}{1}{Énoncé et signification}

\begin{tcolorbox}[thmbox]
  \textbf{Théorème spectral (cas réel).}\enspace%
  \index{théorème spectral}\index{autoadjoint (opérateur)}
  \textit{Soit $E$ un espace euclidien de dimension finie $n$ et
  $f \in \mathcal{L}(E)$ un endomorphisme autoadjoint, i.e.\ vérifiant
  \[
    \forall x,y \in E,\quad \langle f(x),\,y\rangle = \langle x,\,f(y)\rangle.
  \]
  Alors $f$ est diagonalisable dans une base orthonormale. Plus précisément,
  il existe une base orthonormale $(e_1,\ldots,e_n)$ de $E$ et des réels
  $\lambda_1,\ldots,\lambda_n$ tels que $f(e_i) = \lambda_i e_i$ pour tout $i$.}

  \medskip
  \textbf{Version matricielle.}
  Toute matrice symétrique réelle $A \in \mathcal{M}_n(\mathbb{R})$
  (i.e.\ $A^T = A$) est orthogonalement diagonalisable : il existe
  $P \in \mathrm{O}_n(\mathbb{R})$ telle que $P^T A P = \mathrm{diag}(\lambda_1,\ldots,\lambda_n)$
  avec $\lambda_i \in \mathbb{R}$.
\end{tcolorbox}

\begin{significationintuitive}
Un endomorphisme autoadjoint est le pendant abstrait d'une matrice symétrique.
Le théorème spectral affirme que de tels opérateurs sont « géométriquement
simples » : l'espace se décompose en directions propres orthogonales, et
$f$ agit sur chacune par une simple homothétie. C'est le fondement des
méthodes spectrales en physique, probabilités (matrices de covariance)
et traitement du signal (ACP).
\end{significationintuitive}

\decsep{}

%=============================================================
%  § 2 — DÉMONSTRATION
%=============================================================
\secnum{navyblue}{2}{Démonstration}

\begin{ideepreuve}
On montre d'abord que toutes les valeurs propres sont réelles, puis que
les sous-espaces propres associés à des valeurs propres distinctes sont
orthogonaux. On conclut par récurrence sur la dimension : $f$ admet une
valeur propre réelle (son polynôme caractéristique est scindé sur $\mathbb{R}$
par le lemme clé), on se restreint au supplémentaire orthogonal, et on itère.
\end{ideepreuve}

\rubriq{Preuve}

\begin{mdframed}[style=proofbar]
  \textit{Lemme 1 — Valeurs propres réelles.}
  A priori, $\chi_f$ peut admettre des racines complexes ; on travaille donc
  dans le complexifié $E_{\mathbb{C}} = E \otimes_{\mathbb{R}} \mathbb{C}$,
  muni du produit \emph{hermitien} $\langle\cdot,\cdot\rangle$ prolongeant le
  produit scalaire euclidien (sesquilinéaire, semi-linéaire à droite). Comme
  la matrice de $f$ dans une base orthonormale réelle est symétrique réelle,
  donc hermitienne, l'endomorphisme prolongé reste autoadjoint sur
  $E_{\mathbb{C}}$ : $\langle f(x),y\rangle = \langle x,f(y)\rangle$. Soit alors
  $\lambda \in \mathbb{C}$ valeur propre et $v\neq 0$ un vecteur propre associé
  dans $E_{\mathbb{C}}$. Avec $\|v\|^2 = \langle v,v\rangle > 0$ :
  \[
    \lambda\|v\|^2 = \langle f(v),v\rangle = \langle v,f(v)\rangle
    = \overline{\langle f(v),v\rangle} = \bar\lambda\|v\|^2,
  \]
  donc $\lambda = \bar\lambda \in \mathbb{R}$.\hfill$\square$

  \medskip
  \textit{Lemme 2 — Orthogonalité des sous-espaces propres.}
  Si $f(u) = \lambda u$ et $f(v) = \mu v$ avec $\lambda\neq\mu$, alors :
  \[
    \lambda\langle u,v\rangle = \langle f(u),v\rangle = \langle u,f(v)\rangle
    = \mu\langle u,v\rangle,
  \]
  donc $(\lambda-\mu)\langle u,v\rangle = 0$, d'où $\langle u,v\rangle = 0$.\hfill$\square$

  \medskip
  \textit{Lemme 3 — $f$ admet une valeur propre réelle.}
  Le polynôme caractéristique $\chi_f \in \mathbb{R}[X]$ est de degré $n$.
  Par le Lemme 1, toutes ses racines complexes sont réelles, donc $\chi_f$
  est scindé sur $\mathbb{R}$ et admet au moins une racine réelle $\lambda$,
  valeur propre de $f$.\hfill$\square$

  \medskip
  \textit{Récurrence forte sur $\dim E = n$.}

  \textit{Initialisation} : $n=1$, tout endomorphisme est diagonalisable.

  \textit{Hérédité} : on suppose le résultat acquis en toute dimension
  strictement inférieure à $n$. Soit $\lambda$ une valeur propre réelle de $f$ (Lemme 3),
  $V = \ker(f-\lambda\,\mathrm{Id})$ son sous-espace propre, de dimension
  $\dim V \geq 1$ puisque $\lambda$ est valeur propre. On vérifie que
  $V^\perp$ est stable par $f$ : pour $y\in V^\perp$ et $v\in V$,
  \[
    \langle f(y),v\rangle = \langle y,f(v)\rangle = \lambda\langle y,v\rangle = 0.
  \]
  L'espace $V^\perp$, muni de la restriction du produit scalaire de $E$, est
  euclidien ; comme $V^\perp$ est stable par $f$, la restriction $f|_{V^\perp}$
  y est un endomorphisme bien défini, et elle hérite de l'autoadjonction :
  pour $x,y\in V^\perp$, $\langle f(x),y\rangle = \langle x,f(y)\rangle$. Ainsi
  $f|_{V^\perp}$ est autoadjointe sur $V^\perp$, de dimension
  $n - \dim V < n$. Si $V^\perp = \{0\}$ (cas $\dim V = n$, i.e.\
  $f = \lambda\,\mathrm{Id}$), toute base orthonormale de $E$ convient.
  Sinon, par hypothèse de récurrence forte, $V^\perp$ admet une base
  orthonormale de vecteurs propres. En la combinant avec une base orthonormale
  de $V$ (formée de vecteurs propres pour $\lambda$), et comme $E = V\oplus V^\perp$
  avec $V\perp V^\perp$, on obtient une base orthonormale de $E$ diagonalisant $f$.
  \hfill$\blacksquare$
\end{mdframed}

\decsep{}

%=============================================================
%  § 3 — EXEMPLES, CONTRE-EXEMPLES, EXERCICES
%=============================================================
\secnum{exgreen}{3}{Exemples, contre-exemples et exercices}

\rubriq{Exemples}

\begin{mdframed}[style=exbar]
  \textbf{1. Diagonalisation d'une symétrique $2\times 2$.}
  \[
    A = \begin{pmatrix}3 & 1\\1 & 3\end{pmatrix},\quad
    \chi_A = (X-2)(X-4),\quad
    P = \frac{1}{\sqrt{2}}\begin{pmatrix}1&1\\-1&1\end{pmatrix},\quad
    P^TAP = \begin{pmatrix}2&0\\0&4\end{pmatrix}.
  \]

  \smallskip\noindent
  \textbf{2. Analyse en composantes principales (ACP).}
  La matrice de covariance d'un nuage de points est symétrique positive :
  ses vecteurs propres orthonormaux donnent les axes principaux de variance.

  \smallskip\noindent
  \textbf{3. Opérateurs différentiels.}
  L'opérateur $-d^2/dx^2$ sur $L^2([0,\pi])$ avec conditions de Dirichlet
  est autoadjoint ; ses vecteurs propres sont les fonctions $\sin(nx)$,
  base de Fourier en sinus.
\end{mdframed}

\rubriq{Contre-exemples}

\begin{mdframed}[style=cexbar]
  \textbf{La symétrie est indispensable.}\enspace%
  $A = \begin{pmatrix}0&1\\0&0\end{pmatrix}$ n'est pas symétrique :
  $\chi_A = X^2$, la seule valeur propre est $0$ mais $A\neq 0$, donc $A$
  n'est pas diagonalisable.

  \smallskip\noindent
  \textbf{La condition $A^*=A$ ne se relâche pas en \og normale \fg.}\enspace%
  (Ceci n'est pas un contre-exemple au théorème, mais une mise en garde sur
  ses hypothèses sur $\mathbb{C}$.) Le théorème spectral s'étend aux matrices
  hermitiennes ($A^* = A$), dont les valeurs propres sont réelles. Si l'on se
  contente de la normalité ($A^*A = AA^*$), $A$ reste diagonalisable en base
  orthonormale, mais ses valeurs propres ne sont plus réelles en général :
  une matrice unitaire est normale mais en général non hermitienne
  (p.~ex.\ une rotation d'angle $\theta \notin \pi\mathbb{Z}$), et ses valeurs
  propres sont sur le cercle unité. L'hypothèse hermitienne est donc bien ce qui force
  le spectre réel.

  \smallskip\noindent
  \textbf{Dimension infinie.}\enspace%
  En dimension infinie, un opérateur autoadjoint (borné) n'est plus
  nécessairement diagonalisable au sens classique : il faut le calcul
  fonctionnel et la théorie spectrale de von Neumann.
\end{mdframed}

\vspace{6pt}
\exolabel

\begin{mdframed}[style=exobar]
  \begin{enumerate}
    \item Diagonaliser orthogonalement
          $A = \begin{pmatrix}2&1&1\\1&2&1\\1&1&2\end{pmatrix}$
          (valeurs propres $1$, $1$ et $4$)
          et vérifier que $P^TAP$ est diagonale.

    \item Montrer que si $A$ est symétrique et définie positive (i.e.\
          $x^T A x > 0$ pour tout $x\neq 0$), alors toutes ses valeurs
          propres sont strictement positives. Réciproque ?

    \item Soit $f$ autoadjoint sur un espace euclidien. Montrer que
          $\|f\| = \max\{|\lambda|,\;\lambda\text{ v.p. de }f\}$.

    \item (Décomposition polaire) Montrer que toute matrice $M\in\mathcal{M}_n(\mathbb{R})$
          inversible s'écrit $M = QS$ avec $Q$ orthogonale et $S$ symétrique
          définie positive (unicité de $Q$ et de $S$).

    \item Montrer que si $A$ est symétrique réelle, alors $e^A$ est
          symétrique définie positive, et que l'application
          $A\mapsto e^A$ réalise une bijection de l'ensemble des matrices
          symétriques réelles sur l'ensemble des matrices symétriques
          définies positives.
  \end{enumerate}
\end{mdframed}

\leconfooter{D.\ Hilbert (1904) — Théorie spectrale}
